% =========================================================================
\section{Condensadores}\label{sec:condensadores}
% =========================================================================
El sistema de dos conductores en influencia total tiene un nombre propio y una importancia práctica enorme: es el condensador, el dispositivo que
almacena carga (y, como veremos, energía) con una diferencia de potencial dada.

% -------------------------------------------------------------------------
\subsection{Definición y capacidad}\label{sec:condensador-def}
% -------------------------------------------------------------------------
\begin{definicion}[Condensador]\label{def:condensador}
Un \emph{condensador} es un sistema de dos conductores en influencia total cargados con cargas de igual módulo y signo contrario, $Q_1=Q$ y
$Q_2=-Q$ ($Q>0$). Cada uno de los conductores se llama \emph{placa} o \emph{armadura}. Su \emph{capacidad} es el cociente entre el valor absoluto de la carga
común y la diferencia de potencial entre las armaduras:
\begin{equation}\label{eq:cap-condensador}
  \resalta{C=\frac{Q}{\phi_1-\phi_2}=\frac{Q}{\DV} .}
\end{equation}
\end{definicion}

\begin{teorema}[Capacidad en función de los coeficientes]\label{teo:cap-coef}
La capacidad de un condensador solo depende de la geometría de las armaduras y del medio, es positiva, y vale
\begin{equation}\label{eq:cap-coef}
  \resalta{C=\frac{1}{p_{11}+p_{22}-2p_{12}}=\frac{c_{11}c_{22}-c_{12}^{2}}{c_{11}+c_{22}+2c_{12}} .}
\end{equation}
\end{teorema}
\begin{proof}
Haciendo $Q_1=Q$ y $Q_2=-Q$ en las ecuaciones de los coeficientes de potencial \cref{eq:pij},
\[
  \left.\begin{aligned}
    \phi_1&=p_{11}Q_1+p_{12}Q_2=(p_{11}-p_{12})\,Q\\
    \phi_2&=p_{21}Q_1+p_{22}Q_2=(p_{21}-p_{22})\,Q
  \end{aligned}\right\}
  \;\Longrightarrow\;\phi_1-\phi_2=\bigl(p_{11}+p_{22}-2p_{12}\bigr)Q ,
\]
donde hemos usado $p_{21}=p_{12}$. Dividiendo, $C=Q/(\phi_1-\phi_2)=(p_{11}+p_{22}-2p_{12})^{-1}$, que solo depende de la geometría. Es positiva porque
$p_{11}+p_{22}-2p_{12}=\bm{v}\Tt\Pm\,\bm{v}$ con $\bm{v}=(1,-1)\Tt$, y $\Pm$ es definida positiva (\cref{lem:definida}). Para la segunda forma basta sustituir
los elementos de la inversa de $\Cm$: con $D=\det\Cm=c_{11}c_{22}-c_{12}^{2}$, $p_{11}=c_{22}/D$, $p_{22}=c_{11}/D$ y $p_{12}=-c_{12}/D$, de modo que
$p_{11}+p_{22}-2p_{12}=(c_{11}+c_{22}+2c_{12})/D$.
\end{proof}

\begin{corolario}[Conductor encerrado]\label{cor:encerrado}
Si la armadura~1 está encerrada por la armadura~2, entonces
\begin{equation}\label{eq:cap-encerrado}
  C=c_{11}=-c_{12}=-c_{21} .
\end{equation}
En particular, con la armadura exterior a tierra ($\phi_2=0$) y la interior a potencial $V$, se tiene $Q_1=c_{11}V$ y $Q_2=c_{21}V=-Q_1$: es exactamente un
condensador cargado a $\DV=V$, y $C=\tfrac12(c_{11}-c_{21})$.
\end{corolario}
\begin{proof}
Con $\phi_2=0$ y $\phi_1=V$, \cref{eq:cij} da $Q_1=c_{11}V$ y $Q_2=c_{21}V$. Al estar 1 encerrado por 2, todas sus líneas acaban en 2 (influencia total,
\cref{teo:influencia}): $Q_2=-Q_1$, es decir, $c_{21}=-c_{11}$. Entonces $C=Q_1/(\phi_1-\phi_2)=c_{11}$. Restando las dos ecuaciones,
$2Q=(c_{11}-c_{21})V$, que es la última forma. (Esta última expresión \emph{no} es general: solo vale cuando $\phi_2=0$ y hay influencia total.)
\end{proof}

\begin{observacion}[Condensador «generalizado»]\label{obs:generalizado}
La fórmula~\eqref{eq:cap-coef} tiene sentido aunque los dos conductores no estén en influencia total (por ejemplo, dos esferas próximas): es la capacidad
que se mide al transferir una carga $Q$ de un conductor al otro. En un condensador de verdad, en cambio, la influencia total garantiza que el campo
fuera de las armaduras es nulo y que la capacidad no depende de los objetos del entorno.
\end{observacion}

% -------------------------------------------------------------------------
\subsection{Ejemplo: condensador esférico}\label{sec:ej-esferico}
% -------------------------------------------------------------------------
\begin{ejemplo}[Condensador esférico de armaduras delgadas]\label{ej:esferico}
Una esfera conductora de radio $a$ (armadura~1) está rodeada por una cáscara esférica delgada y concéntrica de radio $b>a$ (armadura~2). Calculemos sus
coeficientes y su capacidad.

\emph{Coeficientes de potencial.} Por el \cref{ej:concentricas} (o la regla~\eqref{eq:regla-concentricas}),
\[
  \Pm=\frac{1}{4\pi\eo}\begin{pmatrix}1/a & 1/b\\ 1/b & 1/b\end{pmatrix},\qquad
  C=\frac{1}{p_{11}+p_{22}-2p_{12}}=4\pi\eo\Bigl(\frac1a+\frac1b-\frac2b\Bigr)^{-1}=4\pi\eo\Bigl(\frac1a-\frac1b\Bigr)^{-1},
\]
\begin{equation}\label{eq:cap-esferico}
  \resalta{C=4\pi\eo\,\frac{ab}{b-a} .}
\end{equation}

\emph{Coeficientes de capacidad e influencia.} El determinante es
\[
  \det\Pm=\Bigl(\frac{1}{4\pi\eo}\Bigr)^{2}\Bigl(\frac1a\,\frac1b-\frac1b\,\frac1b\Bigr)=\Bigl(\frac{1}{4\pi\eo}\Bigr)^{2}\frac{b-a}{ab^{2}} ,
\]
y la inversa de una matriz $2\times2$ es $\Pm^{-1}=\frac{1}{\det\Pm}\bigl(\begin{smallmatrix}p_{22}&-p_{12}\\-p_{12}&p_{11}\end{smallmatrix}\bigr)$:
\[
  \Cm=(4\pi\eo)^{2}\,\frac{ab^{2}}{b-a}\,\frac{1}{4\pi\eo}\begin{pmatrix}1/b & -1/b\\ -1/b & 1/a\end{pmatrix}
  =4\pi\eo\,\frac{ab}{b-a}\begin{pmatrix}1 & -1\\ -1 & b/a\end{pmatrix}.
\]
Se comprueban todas las propiedades: $\Cm$ es simétrica, $c_{11},c_{22}>0$ y $c_{12}<0$. Además, $c_{11}=-c_{21}=C$: la esfera interior está en influencia
total (\cref{cor:encerrado}). En cambio, $c_{22}+c_{12}=4\pi\eo\,\frac{b^{2}-ab}{b-a}=4\pi\eo\,b>0$: parte de las líneas que salen de la cáscara se van al
infinito, y lo que «sobra» es justamente la capacidad de una esfera aislada de radio $b$.

\emph{Placas próximas.} Si $b=a+d$ con $d\ll a$,
$C=4\pi\eo\,\frac{a(a+d)}{d}\simeq4\pi\eo\,\frac{a^{2}}{d}=\eo\,\frac{A}{d}$ con $A=4\pi a^{2}$: el condensador esférico se comporta como uno plano de
área $A$ y separación $d$ (\cref{sem:plano}).
\end{ejemplo}

% =========================================================================
%  Figura: condensador esférico (b = 2a).
%  (a) Sección: +Q en la esfera, -Q en la cara interior de la cáscara, campo
%      radial entre armaduras y gaussiana a < r < b.
%  (b) E/E0 frente a r/a, con E0 = kQ/a^2: el área sombreada es DeltaV/phi0.
%  (c) phi/phi0 con la armadura exterior a tierra, phi0 = kQ/a.
% =========================================================================
\begin{figure}[t]
\centering
\begin{tikzpicture}[font=\small]
  % ---------------------------------------------------------------- (a)
  \begin{scope}[xshift=7.6cm, x=0.85cm, y=0.85cm]
    \foreach \a in {15,45,...,345}{
      \draw[azulacero, line width=0.55pt, -{Stealth[length=3.5pt]}] (\a:1.08) -- (\a:1.6);
      \draw[azulacero, line width=0.55pt] (\a:1.6) -- (\a:1.9);}
    \draw[azulprof, line width=1pt] (0,0) circle (2);
    \shade[ball color=azulvelo] (0,0) circle (1);
    \draw[azulprof, line width=0.6pt] (0,0) circle (1);
    \foreach \a in {0,45,...,315}{\node[signopos, font=\tiny\bfseries] at (\a:0.82) {$+$};
                                 \node[signoneg, font=\tiny\bfseries] at (\a+22.5:1.86) {$-$};}
    \draw[gauss] (0,0) circle (1.45);
    \node[font=\footnotesize, text=cargapos] at (0,0.3) {$+Q$};
    \node[font=\footnotesize] at (0,-0.25) {$a$};
    \node[font=\footnotesize, text=carganeg] at (2.25,1.55) {$-Q$};
    \node[font=\footnotesize] at (2.05,-1.5) {$b$};
    \node[font=\footnotesize, text=grispizarra, align=center] at (0,-2.55) {$\vE=\vc{0}$ fuera de $b$};
    \node at (0,-3.2) {(a)};
  \end{scope}
  % ---------------------------------------------------------------- (b)
  \begin{scope}[xshift=1.2cm, yshift=-7.0cm]
    \begin{axis}[tema, width=6.6cm, height=4.6cm, xmin=0, xmax=3.2, ymin=0, ymax=1.15,
        xlabel={$r/a$}, ylabel={$E/E_0$}, xtick={0,1,2,3}, ytick={0,0.25,0.5,1}]
      \addplot[draw=none, fill=azulvelo, domain=1:2, samples=40] {1/x^2} \closedcycle;
      \addplot[azulprof, line width=1.1pt, domain=1:2, samples=40] {1/x^2};
      \addplot[azulprof, line width=1.1pt] coordinates {(0,0) (1,0)};
      \addplot[azulprof, line width=1.1pt] coordinates {(2,0) (3.2,0)};
      \draw[punteado] (axis cs:1,0) -- (axis cs:1,1) (axis cs:2,0) -- (axis cs:2,0.25);
      \node[font=\scriptsize] at (axis cs:1.45,0.18) {$\DV$};
      \node[font=\scriptsize, text=grispizarra] at (axis cs:0.5,0.12) {$E=0$};
      \node[font=\scriptsize, text=grispizarra] at (axis cs:2.6,0.12) {$E=0$};
    \end{axis}
    \node at (2.6,-1.35) {(b)};
  \end{scope}
  % ---------------------------------------------------------------- (c)
  \begin{scope}[xshift=9.0cm, yshift=-7.0cm]
    \begin{axis}[tema, width=6.6cm, height=4.6cm, xmin=0, xmax=3.2, ymin=0, ymax=0.62,
        xlabel={$r/a$}, ylabel={$\phi/\phi_0$}, xtick={0,1,2,3}, ytick={0,0.25,0.5}]
      \addplot[azulacero, line width=1.1pt] coordinates {(0,0.5) (1,0.5)};
      \addplot[azulacero, line width=1.1pt, domain=1:2, samples=40] {1/x-0.5};
      \addplot[azulacero, line width=1.1pt] coordinates {(2,0) (3.2,0)};
      \draw[punteado] (axis cs:1,0) -- (axis cs:1,0.5) (axis cs:2,0) -- (axis cs:2,0.6);
      \draw[{Stealth[length=3pt]}-{Stealth[length=3pt]}, cargapos] (axis cs:2.35,0) -- (axis cs:2.35,0.5);
      \draw[punteado] (axis cs:1,0.5) -- (axis cs:2.45,0.5);
      \node[font=\scriptsize, text=cargapos, anchor=west] at (axis cs:2.4,0.27) {$\DV$};
    \end{axis}
    \node at (2.6,-1.35) {(c)};
  \end{scope}
\end{tikzpicture}
\caption{Condensador esférico con $b=2a$. (a) Sección: toda la carga $-Q$ de la armadura exterior está en su cara interior, y la gaussiana entre armaduras
solo encierra $+Q$. (b) Campo, con $E_0=\kc Q/a^{2}$: es nulo dentro de la esfera y fuera de la cáscara, y el área bajo la curva entre $a$ y $b$ es
$\DV/\phi_0$. (c) Potencial con la armadura exterior a tierra, $\phi_0=\kc Q/a$: constante en la esfera y nulo desde la cáscara hacia fuera.}
\label{fig:condensador-esferico}
\end{figure}


\begin{ejemplo}[El mismo sistema, en otras situaciones]\label{ej:esferico-situaciones}
Los coeficientes resuelven cualquier situación del sistema, no solo la del condensador.

\emph{Esfera interior a tierra.} Si $\phi_1=0$ y $\phi_2\neq0$, $Q_2=c_{22}\phi_2$ y $Q_1=c_{12}\phi_2$, de modo que
\[
  \frac{Q_1}{Q_2}=\frac{c_{12}}{c_{22}}=-\frac{a}{b} :
\]
el sistema \emph{no} está en influencia total respecto de la cáscara, y la esfera interior adquiere menos carga (en valor absoluto) que la exterior. Si
$b=2a$, $Q_1=-Q_2/2$: la cara interior de la cáscara lleva $+Q_2/2$ y la exterior $+Q_2/2$.

\emph{Cargas dadas.} Si conocemos $Q_1$ y $Q_2$, los potenciales salen directamente con los coeficientes de potencial,
\[
  \phi_1=p_{11}Q_1+p_{12}Q_2=\frac{1}{4\pi\eo}\Bigl(\frac{Q_1}{a}+\frac{Q_2}{b}\Bigr),\qquad
  \phi_2=p_{21}Q_1+p_{22}Q_2=\frac{1}{4\pi\eo}\,\frac{Q_1+Q_2}{b} ,
\]
mientras que con los coeficientes de capacidad habría que resolver un sistema lineal. Es útil comprobarlo resolviendo el problema «como en el Tema~1»:
por Gauss,
\[
  \vE(\vc{r})=\begin{cases}\dfrac{Q_1+Q_2}{4\pi\eo r^{2}}\,\ur, & r>b,\\[8pt] \dfrac{Q_1}{4\pi\eo r^{2}}\,\ur, & a<r<b,\\[8pt] \vc{0}, & r<a,\end{cases}
\]
con un salto de $\vE$ en $r=b$ debido a la carga de la cáscara. El potencial es continuo, con $\phi(\infty)=0$: para $r\ge b$,
$\phi=\frac{Q_1+Q_2}{4\pi\eo r}$; para $a\le r\le b$,
\[
  \phi(r)-\phi(b)=-\int_{b}^{r}\vE\cdot\dif\vc{r}'=\int_{r}^{b}\frac{Q_1}{4\pi\eo r'^{2}}\,\dif r'=\frac{Q_1}{4\pi\eo}\Bigl(\frac1r-\frac1b\Bigr)
  \;\Longrightarrow\;\phi(r)=\frac{1}{4\pi\eo}\Bigl(\frac{Q_1}{r}+\frac{Q_2}{b}\Bigr),
\]
y para $r\le a$, $\phi=\frac{1}{4\pi\eo}\bigl(\frac{Q_1}{a}+\frac{Q_2}{b}\bigr)$, que coincide con $\phi_1$. Con los coeficientes de potencial el cálculo es inmediato; con los de
capacidad, para pasar de cargas a potenciales, habría que invertir la matriz.
\end{ejemplo}

% -------------------------------------------------------------------------
\subsection{Ejemplo: condensador esférico con armadura exterior gruesa}\label{sec:ej-grueso}
% -------------------------------------------------------------------------
\begin{ejemplo}[Esfera y corteza esférica gruesa]\label{ej:grueso}
El conductor~1 es una esfera de radio $a$; el conductor~2, una corteza esférica concéntrica de radio interior $b$ y exterior $c$. Escribimos
$\kc=1/(4\pi\eo)$.

\emph{Caso A: $Q_1=Q$, $Q_2=0$.} Por influencia total, la cara interior de la corteza adquiere $-Q$ (Gauss con una superficie dentro del metal) y, por
conservación de la carga, la exterior $+Q$. Fuera de la corteza el campo es el de una carga $Q$ en el centro, y $\phi(\vc{r})=-\int_{\infty}^{\vc{r}}\vE\cdot\dl$
da, sobre la corteza y sobre la esfera,
\[
  \phi_2=\int_{c}^{\infty}\frac{\kc Q}{r^{2}}\,\dif r=\frac{\kc Q}{c},\qquad
  \phi_1=\phi_2+\int_{a}^{b}\frac{\kc Q}{r^{2}}\,\dif r=\frac{\kc Q}{c}+\kc Q\Bigl(\frac1a-\frac1b\Bigr),
\]
luego $p_{11}=\kc\bigl(\frac1a-\frac1b+\frac1c\bigr)$ y $p_{21}=\kc/c$.

\emph{Caso B: $Q_1=0$, $Q_2=Q$.} La carga se reparte en la cara exterior de la corteza; dentro el campo es nulo y la esfera interior queda al mismo potencial
que la corteza, $\phi_1=\phi_2=\kc Q/c$. Así $p_{12}=p_{22}=\kc/c$, y
\[
  \Pm=\frac{1}{4\pi\eo}\begin{pmatrix}\dfrac1a-\dfrac1b+\dfrac1c & \dfrac1c\\[8pt] \dfrac1c & \dfrac1c\end{pmatrix}.
\]

\emph{Capacidad del condensador.}
\[
  C=\frac{Q}{\phi_1-\phi_2}=\frac{1}{p_{11}+p_{22}-2p_{12}}=\frac{4\pi\eo}{\frac1a-\frac1b+\frac1c+\frac1c-\frac2c}=4\pi\eo\,\frac{ab}{b-a} ,
\]
la misma que con armaduras delgadas: el espesor de la armadura exterior no influye, porque entre las armaduras el campo solo depende de $Q$. Con $d=b-a$,
$C=\eo\,A_{\mathrm{ef}}/d$ con $A_{\mathrm{ef}}=4\pi ab=\sqrt{4\pi a^{2}\cdot4\pi b^{2}}$, la media geométrica de las áreas de las armaduras.

\emph{Inversión de la matriz.} Con $\det\Pm=p_{11}p_{22}-p_{12}^{2}=\kc^{2}\,\frac{b-a}{abc}$ y la matriz adjunta,
$\Pm^{-1}=\frac{1}{\det\Pm}\,[\adj\Pm]\Tt=\frac{1}{\det\Pm}\bigl(\begin{smallmatrix}p_{22}&-p_{12}\\-p_{21}&p_{11}\end{smallmatrix}\bigr)$. Como $\Pm$ es simétrica, la
matriz de cofactores también lo es, y trasponer no cambia nada: por eso $\Cm$ sale simétrica automáticamente. Resulta
\[
  c_{11}=\frac{p_{22}}{\det\Pm}=4\pi\eo\,\frac{ab}{b-a},\qquad c_{12}=c_{21}=-\frac{p_{12}}{\det\Pm}=-4\pi\eo\,\frac{ab}{b-a},
\]
\[
  c_{22}=\frac{p_{11}}{\det\Pm}=4\pi\eo\Bigl(c+\frac{ab}{b-a}\Bigr).
\]
\emph{Interpretación.} $c_{11}>0$ y $c_{12}<0$; $c_{11}=-c_{21}$ es la igualdad de la influencia total, pues todas las líneas de la esfera interior acaban en la
corteza; en cambio $c_{22}>-c_{12}$, porque parte de las líneas de la corteza van al infinito: $c_{22}=C+4\pi\eo c$, la capacidad del condensador más la de
una esfera aislada de radio $c$.
\end{ejemplo}

% =========================================================================
%  Figura: esfera (a) y corteza gruesa (b, c), con a = 1, b = 2, c = 2,5.
%  Potencial normalizado a kQ/a en los dos estados elementales:
%  Caso A (Q1 = Q, Q2 = 0): 1/x (x >= c); 1/c (b <= x <= c);
%                           1/c + 1/x - 1/b (a <= x <= b); 1/c + 1/a - 1/b (x <= a)
%  Caso B (Q1 = 0, Q2 = Q): 1/x (x >= c); 1/c (x <= c)
% =========================================================================
\begin{figure}[t]
\centering
\begin{tikzpicture}[font=\small]
  % ---------------------------------------------------------------- (a)
  \begin{scope}[x=0.72cm, y=0.72cm]
    \fill[metal!65!white, even odd rule] (0,0) circle (2.5) (0,0) circle (2);
    \fill[pattern=north east lines, pattern color=azulcielo!70, even odd rule] (0,0) circle (2.5) (0,0) circle (2);
    \draw[azulprof, line width=0.6pt] (0,0) circle (2.5) (0,0) circle (2);
    \shade[ball color=azulvelo] (0,0) circle (1);
    \draw[azulprof, line width=0.6pt] (0,0) circle (1);
    \draw[-{Stealth[length=3.5pt]}, azulprof, line width=0.5pt] (0,0) -- (90:1) node[pos=0.5, right=-1pt, font=\scriptsize] {$a$};
    \draw[-{Stealth[length=3.5pt]}, azulprof, line width=0.5pt] (0,0) -- (210:2) node[pos=0.6, above left=-2pt, font=\scriptsize] {$b$};
    \draw[-{Stealth[length=3.5pt]}, azulprof, line width=0.5pt] (0,0) -- (-30:2.5) node[pos=0.8, above right=-2pt, font=\scriptsize] {$c$};
    \node[font=\footnotesize] at (0,-0.35) {$1$};
    \node[font=\footnotesize, fill=white, inner sep=1pt] at (0,2.25) {$2$};
    \node at (0,-3.45) {(a)};
  \end{scope}
  % ---------------------------------------------------------------- (b)
  \begin{scope}[xshift=4.7cm, yshift=-2.1cm]
    \begin{axis}[tema, width=6.4cm, height=4.6cm, xmin=0, xmax=4.5, ymin=0, ymax=1.02,
        xlabel={$r/a$}, ylabel={$\phi\,/\,(\kc Q/a)$}, xtick={0,1,2,2.5,4}, xticklabels={$0$,$1$,$2$,$2{,}5$,$4$},
        ytick={0,0.4,0.9}, legend style={at={(0.98,0.97)}, anchor=north east}]
      \fill[metal!45!white] (axis cs:2,0) rectangle (axis cs:2.5,1.02);
      \addplot[azulprof, line width=1.1pt] coordinates {(0,0.9) (1,0.9)};
      \addlegendentry{caso A: $Q_1=Q$}
      \addplot[azulprof, line width=1.1pt, domain=1:2, samples=40, forget plot] {0.4+1/x-0.5};
      \addplot[azulprof, line width=1.1pt, forget plot] coordinates {(2,0.4) (2.5,0.4)};
      \addplot[azulprof, line width=1.1pt, domain=2.5:4.5, samples=40, forget plot] {1/x};
      \addplot[azulcielo, line width=1.1pt, densely dashed] coordinates {(0,0.4) (2.5,0.4)};
      \addlegendentry{caso B: $Q_2=Q$}
      \addplot[azulcielo, line width=1.1pt, densely dashed, domain=2.5:4.5, samples=40, forget plot] {1/x};
      \draw[{Stealth[length=3pt]}-{Stealth[length=3pt]}, cargapos] (axis cs:0.5,0.4) -- (axis cs:0.5,0.9)
        node[midway, right, font=\scriptsize] {$\phi_1-\phi_2$};
    \end{axis}
    \node at (2.45,-1.35) {(b)};
  \end{scope}
\end{tikzpicture}
\caption{Condensador esférico con armadura exterior gruesa ($a=1$, $b=2$, $c=2{,}5$). (a) Geometría: la corteza, entre $b$ y $c$, es un volumen
equipotencial. (b) Potencial en los dos estados que definen los coeficientes de potencial: con la esfera cargada (caso A) el potencial cae entre $a$ y $b$, y
la diferencia $\phi_1-\phi_2$ solo depende de $a$ y $b$; con la corteza cargada (caso B) el interior queda a potencial constante $\kc Q/c$.}
\label{fig:corteza-gruesa}
\end{figure}


% -------------------------------------------------------------------------
\subsection{Los tres condensadores canónicos}\label{sec:canonicos}
% -------------------------------------------------------------------------
Para geometrías indefinidas (cilindros o planos infinitos) no se puede tomar $\phi(\infty)=0$, porque el potencial de una distribución infinita de carga
diverge; los coeficientes de potencial solo pueden definirse respecto de una referencia arbitraria (\cref{sem:cilindrico}). En un condensador, sin
embargo, esto no supone ningún problema: el campo fuera de las armaduras es nulo y solo interesa la diferencia de potencial, que es finita. La
\cref{tab:canonicos} reúne los tres casos clásicos, que se deducen en detalle en los seminarios (\cref{sem:esferico,sem:cilindrico,sem:plano}).

\begin{table}[ht]
\centering
\caption{Los tres condensadores canónicos. El campo entre las armaduras se obtiene con la ley de Gauss, la diferencia de potencial integrándolo de la
armadura positiva a la negativa, y la capacidad como $C=Q/\DV$.}
\label{tab:canonicos}
\begin{tabular}{@{}lccc@{}}
  \toprule
  Condensador & Campo entre armaduras & $\DV$ & Capacidad\\
  \midrule
  Esférico (radios $a<b$) & $\dfrac{Q}{4\pi\eo r^{2}}$ & $\dfrac{Q}{4\pi\eo}\Bigl(\dfrac1a-\dfrac1b\Bigr)$ & $4\pi\eo\,\dfrac{ab}{b-a}$\\[12pt]
  Cilíndrico (radios $a<b$, longitud $\ell$) & $\dfrac{\lambda}{2\pi\eo r}$ & $\dfrac{\lambda}{2\pi\eo}\ln\dfrac{b}{a}$ & $\dfrac{2\pi\eo\,\ell}{\ln(b/a)}$\\[12pt]
  Plano (área $A$, separación $d$) & $\dfrac{Q}{\eo A}$ & $\dfrac{Qd}{\eo A}$ & $\eo\,\dfrac{A}{d}$\\[8pt]
  \bottomrule
\end{tabular}
\end{table}

\begin{observacion}[Un límite común]\label{obs:limite-plano}
Cuando la separación entre armaduras es pequeña frente a sus radios de curvatura, los tres casos tienden a la misma expresión $C\simeq\eo A/d$: localmente,
dos superficies muy próximas se ven planas. Para el cilíndrico, con $b=a+d$, $\ln(1+d/a)\simeq d/a$ y $C/\ell\simeq2\pi\eo a/d=\eo(A/\ell)/d$.
\end{observacion}
